\documentclass[11pt,a4paper]{article}

\usepackage[T1]{fontenc}
\usepackage[utf8]{inputenc}
\usepackage[polish]{babel}
\usepackage{lmodern}
\usepackage[margin=2.2cm]{geometry}
\usepackage{amsmath, amssymb}
\usepackage{mathtools}
\usepackage{booktabs}
\usepackage{array}
\usepackage{enumitem}
\usepackage{xcolor}
\usepackage[hidelinks]{hyperref}
\usepackage{titlesec}
\usepackage{tcolorbox}
\tcbuselibrary{breakable, skins}

% --- styl -----------------------------------------------------------------
\definecolor{accent}{HTML}{4F46E5}
\definecolor{soft}{HTML}{F3F4F6}
\definecolor{rule}{HTML}{D1D5DB}

\titleformat{\section}{\Large\bfseries\color{accent}}{\thesection}{0.6em}{}
\titleformat{\subsection}{\large\bfseries}{\thesubsection}{0.6em}{}
\setlist[itemize]{itemsep=2pt, topsep=2pt, leftmargin=1.4em}
\setlist[description]{itemsep=2pt, topsep=2pt, leftmargin=1.4em}

\newtcolorbox{paformula}[1][]{
  colback=soft, colframe=accent, boxrule=0.4pt,
  arc=2pt, left=10pt, right=10pt, top=8pt, bottom=8pt,
  fonttitle=\bfseries,
  #1
}

\title{\textbf{Metryki Psychoakustyczne dla Analizy Hałasu Dronów}\\
       \large Raport matematyczny: SQM i Uciążliwość Psychoakustyczna (PA)}
\author{Źródło: Vaiuso, Righi, Coretti, Apicella (ZHAW, 2024) ---
        \href{https://arxiv.org/abs/2410.22208}{arXiv:2410.22208}}
\date{Skompilowano \today}

\begin{document}
\maketitle

\begin{abstract}
\noindent
Dokument zawiera podstawy matematyczne potrzebne do dodania zakładki
\textbf{Psychoakustyka} w SoundVisualizer. Obejmuje cztery metryki jakości
dźwięku (SQM, \emph{Sound Quality Metrics}) --- \emph{głośność}, \emph{ostrość},
\emph{chropowatość}, \emph{siłę fluktuacji} --- oraz złożony wskaźnik
\emph{uciążliwości psychoakustycznej} (PA, \emph{Psychoacoustic Annoyance}),
który łączy je w jedną liczbę. Same wzory pochodzą z modelu Zwickera i Fastla;
przywoływany artykuł stosuje je do hałasu dronów. Wszystkie metryki opierają
się na modelu \textbf{pasm krytycznych Bark} (24 pasma w zakresie słyszalnym).
\end{abstract}

% --- 1. Głośność ----------------------------------------------------------
\section{Głośność \texorpdfstring{$L$}{L} (son)}

\textit{Jak głośny jest dźwięk w odczuciu słuchacza.} Odniesienie:
$1\ \text{son} = $ ton sinusoidalny 1\,kHz o poziomie 40\,dB SPL. Norma:
ISO 532-1 / DIN 45631 (metoda Zwickera).

\[
  L = \int_{z=0}^{24} L'(z)\, \mathrm{d}z
\]

\begin{description}
  \item[$L'(z)$] \emph{głośność właściwa} jako funkcja krytycznego pasma $z$ w
                jednostkach Bark (jednostka: $\text{son}/\text{Bark}$).
                Obliczana w każdym paśmie z poziomu SPL przez nieliniowe
                prawo potęgowe z korekcjami maskowania międzypasmowego
                (algorytm ISO 532B).
  \item[$z$]    częstotliwość krytyczna pasma, $0 \le z \le 24$\,Bark.
                Odwzorowanie na skalę częstotliwości jest nieliniowe:
                $1$\,Bark~$\approx 100$\,Hz w zakresie niskich częstotliwości,
                rośnie do $\approx 20$\,kHz przy $24$\,Bark.
\end{description}

\noindent\textbf{Dwa warianty} są dla nas istotne:
\begin{itemize}
  \item \textbf{Głośność stacjonarna} (ISO 532B / DIN 45631) --- pojedyncza
        wartość dla całego nagrania.
  \item \textbf{Głośność zmienna w czasie} $L(t)$ (DIN 45631/A) --- liczona
        w odstępach $\sim 2$\,ms. Niezbędna do wyznaczenia PA, gdzie
        wykorzystywany jest 5-ty percentyl $N_5$ (głośność przekroczona przez
        5\,\% czasu trwania --- ujmuje krótkie szczyty, a nie poziom średni).
\end{itemize}

% --- 2. Ostrość -----------------------------------------------------------
\section{Ostrość \texorpdfstring{$S$}{S} (acum)}

\textit{Jak ostry, przeszywający odbierany jest dźwięk.} Odniesienie:
$1\ \text{acum} = $ szum wąskopasmowy wycentrowany w 1\,kHz, poziom 60\,dB,
szerokość pasma 160\,Hz. Artykuł stosuje wagi z DIN 45692 (rozwinięcie metody
Auresa / Von Bismarcka).

\[
  S = \frac{k}{L} \int_{z=0}^{24} L'(z)\, g(z)\, z\, \mathrm{d}z
\]

\begin{description}
  \item[$g(z)$] funkcja wagowa wzmacniająca wkład wysokich pasm krytycznych
                (wysokie wartości $z$).
  \item[$k$]    stała skalująca zależna od metody; dobierana tak, aby sygnał
                odniesienia dawał dokładnie $1$\,acum.
  \item[$L$]    całkowita głośność z poprzedniej sekcji.
\end{description}

\noindent
Interpretacja geometryczna: $S$ to \emph{ważony środek ciężkości głośności
właściwej na osi Bark}, przesunięty w stronę wysokich częstotliwości przez
$g(z)$. Dźwięki z większą energią w wysokich pasmach mają wyższą ostrość.

% --- 3. Chropowatość ------------------------------------------------------
\section{Chropowatość \texorpdfstring{$R$}{R} (asper)}

\textit{Jak chropowaty, agresywny w odbiorze jest dźwięk.} Powstaje wskutek
modulacji amplitudy w zakresie $30$--$300$\,Hz --- szybkich fluktuacji, których
ucho nie potrafi rozróżnić jako odrębnych zdarzeń. Odniesienie:
$1\ \text{asper} = $ ton 60\,dB / 1\,kHz, w 100\,\% modulowany sinusoidą o
częstotliwości 70\,Hz.

\[
  R = \mathrm{cal} \cdot \int_{z=0}^{24} f_{\text{mod}}\, \Delta L\, \mathrm{d}z
\]

\begin{description}
  \item[$f_{\text{mod}}$] częstotliwość modulacji w każdym paśmie krytycznym
                          (Hz), wyznaczana przez wykrywanie szczytów w widmie
                          obwiedni sygnału.
  \item[$\Delta L$]       głębokość modulacji w paśmie:
                          $\max(\text{env}) - \min(\text{env})$ w dB.
  \item[cal]              stała kalibracyjna gwarantująca, że sygnał
                          odniesienia daje $1$\,asper.
\end{description}

\noindent
Hałas dronów zazwyczaj wykazuje \emph{wysoką} chropowatość: częstotliwość
przejścia łopaty (BPF) i jej harmoniczne wzajemnie się modulują, dając
obwiednię w istotnym akustycznie zakresie 30--300\,Hz.

% --- 4. Siła fluktuacji ---------------------------------------------------
\section{Siła fluktuacji \texorpdfstring{$F$}{F} (vacil)}

\textit{Jak "falujący" jest dźwięk.} Wynika z \emph{wolnej} modulacji
amplitudy: $0$--$30$\,Hz, ze szczytem czułości przy $4$\,Hz. Odniesienie:
$1\ \text{vacil} = $ ton 60\,dB / 1\,kHz, w 100\,\% modulowany sinusoidą
o częstotliwości 4\,Hz.

\[
  F = \frac{0.008 \cdot \displaystyle\int_{z=0}^{24} \Delta L\, \mathrm{d}z}{\dfrac{f_{\text{mod}}}{4} + \dfrac{4}{f_{\text{mod}}}}
\]

\noindent
Struktura wzoru jest taka sama jak dla chropowatości, ale mianownik
$f_{\text{mod}}/4 + 4/f_{\text{mod}}$ osiąga \emph{minimum} (a $F$ swoje
maksimum) w $f_{\text{mod}} = 4$\,Hz --- odpowiada to empirycznemu szczytowi
czułości ucha na powolne fluktuacje.

Dla dronów wartość $F$ jest typowo bliska zeru ($\sim 0.01$\,vacil w pomiarach
zawisu w przywoływanym artykule) --- modulacje silnika wirnikowego leżą wyraźnie
powyżej zakresu wolnych fluktuacji.

% --- 5. PA ----------------------------------------------------------------
\section{Uciążliwość psychoakustyczna \texorpdfstring{$PA$}{PA}}

Złożony wskaźnik łączący cztery powyższe metryki w jedną liczbę. Postać
Zwickera i Fastla:

\begin{paformula}
\[
  PA = N_5 \left(1 + \sqrt{w_S^2 + w_{FR}^2}\right)
\]
\end{paformula}

\noindent
Szczytowa głośność $N_5$ jest mnożona przez bezwymiarowy "współczynnik
nieprzyjemności" $(1 + \sqrt{w_S^2 + w_{FR}^2})$. Dwie wagi:

\paragraph{Waga ostrości \boldmath$w_S$:}
\[
  w_S =
  \begin{cases}
    (S - 1.75) \cdot 0.25 \cdot \log_{10}(N_5 + 10), & \text{jeśli } S > 1.75 \\
    0, & \text{jeśli } S \le 1.75
  \end{cases}
\]
Poniżej progu 1.75\,acum ostrość nie wnosi wkładu. Powyżej progu wkład rośnie
wraz z nadwyżką ostrości i całkowitą głośnością.

\paragraph{Łączna waga chropowatości i fluktuacji \boldmath$w_{FR}$:}
\[
  w_{FR} = \frac{2.18}{N_5^{0.4}}\,(0.4\,F + 0.6\,R)
\]
Chropowatość waży 1{,}5$\times$ więcej niż siła fluktuacji. Czynnik
$N_5^{-0.4}$ łagodzi wpływ tych składowych przy już głośnych dźwiękach.

\paragraph{Model intuicyjny.}
$PA \approx \text{głośność} \times \text{mnożnik nieprzyjemności barwy}$.
Głośność dominuje; ostrość, chropowatość i fluktuacja ją modulują.

% --- 6. Wartości typowe ---------------------------------------------------
\section{Obserwowane wartości dla dronów}

Artykuł podaje następujące wartości dla DJI Matrice 300 RTK w fazie zawisu:

\begin{center}
\begin{tabular}{lrrl}
  \toprule
  Metryka & Średnia & Maks. & Jednostka \\
  \midrule
  Głośność ($L$)         & 19{,}30 & 29{,}56 & son    \\
  Ostrość ($S$)          & 1{,}85  & 3{,}69  & acum   \\
  Chropowatość ($R$)     & 0{,}28  & 0{,}95  & asper  \\
  Fluktuacja ($F$)       & 0{,}01  & 0{,}06  & vacil  \\
  Uciążliwość ($PA$)     & 20{,}31 & 32{,}57 & ---    \\
  \bottomrule
\end{tabular}
\end{center}

\noindent
Wnioski:
\begin{itemize}
  \item Ostrość leży \emph{tuż} powyżej progu aktywacji 1{,}75\,acum ---
        wkład $w_S$ jest niewielki.
  \item Siła fluktuacji praktycznie nie wpływa na $w_{FR}$.
  \item Chropowatość jest dominującą składową barwy poprzez człon
        $0{,}6 \cdot R$.
  \item Wskutek tego $PA$ dla drona w zawisie wynosi około $L \times 1{,}05$
        --- składowe barwowe podnoszą wskaźnik tylko o kilka procent ponad
        samą głośność.
\end{itemize}

% --- 7. Implementacja -----------------------------------------------------
\section{Ścieżka implementacji}

W artykule obliczenia wykonano w MATLAB Audio Toolbox (oprogramowanie
własnościowe). Otwartoźródłowe alternatywy w Pythonie:

\begin{center}
\renewcommand{\arraystretch}{1.2}
\begin{tabular}{>{\bfseries}l p{6.5cm} p{4.5cm}}
  \toprule
  Biblioteka & Co udostępnia & Wady i zalety \\
  \midrule
  \texttt{mosqito} &
    Wszystkie cztery SQM oraz PA, zgodne z ISO/DIN, licencja BSD,
    instalacja przez pip &
    Dodaje zależność ($\sim 10$\,MB), ale jest naturalnym wyborem \\

  \texttt{scikit-maad}, \texttt{pyAudioAnalysis} &
    Głośność i część SQM &
    Mniej kompletne; PA brak \\

  Implementacja własna &
    Pełna kontrola i zrozumienie &
    Tygodnie pracy; nietrywialne filtry i model maskowania
    ISO 532B \\
  \bottomrule
\end{tabular}
\end{center}

\noindent
\textbf{Rekomendacja:} dodać \texttt{mosqito} do \texttt{pyproject.toml}.
Po stronie serwera obliczać pięć wartości raz na pomiar akustyczny i
buforować je obok \texttt{meta.json}. Frontend tylko odczytuje liczby ---
brak obliczeń w przeglądarce.

\paragraph{Cytowanie.}
Jeżeli publikujemy wyniki uzyskane z użyciem MOSQITO, cytujemy je w sposób
wskazany przez autorów:
\begin{quote}
\small
Green Forge Coop. \textit{MOSQITO} [Computer software].
\href{https://doi.org/10.5281/zenodo.5284054}{https://doi.org/10.5281/zenodo.5284054}
\end{quote}
(Konkretne wydanie najwygodniej zacytować przyciskiem \emph{Cite this repository}
na stronie projektu MOSQITO w GitHubie.)

\paragraph{Proponowane widoki zakładki Psychoakustyka.}
\begin{enumerate}
  \item \textbf{Tabela per-mikrofon dla wybranego punktu PWM} --- jeden wiersz
        na mikrofon, kolumny $L$, $S$, $R$, $F$, $PA$. Ten sam format co
        nagłówek wydajnościowy w zakładce FFT, ale w rozbiciu na mikrofony.
  \item \textbf{Metryki w funkcji elewacji} --- małe wykresy liniowe dla
        każdej metryki w funkcji kąta elewacji. Ujawniają przestrzenny
        wzorzec składowych nieprzyjemności, a nie tylko głośności.
  \item \textbf{Metryki w funkcji PWM (lub RPM, ciągu)} --- analogicznie do
        wykresu rozrzutu zakładki Custom, z dowolną metryką SQM lub $PA$
        na osi Y. Odpowiada na pytania w stylu: „przy jakich obrotach prop
        zaczyna brzmieć uciążliwie nie z powodu głośności?"
\end{enumerate}

\paragraph{Otwarte kwestie na etapie implementacji.}
\begin{itemize}
  \item \textbf{Kalibracja.} Jednostki son / acum / asper / vacil zakładają,
        że sygnał wejściowy jest wyrażony w dB SPL, nie w dBFS. Wymagane
        jest zastosowanie kalibracji UMIK-2 \emph{przed} przekazaniem audio
        do \texttt{mosqito}. Działa tak, jeśli dla mikrofonu wgrano plik
        kalibracyjny; w przeciwnym razie wartości bezwzględne będą
        zafałszowane (porównania względne pozostają poprawne).
  \item \textbf{Długość okna.} Dla głośności zmiennej w czasie $\to N_5$
        potrzebujemy $\gtrsim 1$\,s nagrania. Nasze typowe okno akwizycji
        to 1{,}5\,s na krok PWM --- marginalnie wystarczające.
  \item \textbf{Porównanie z Norsonic.} Po dostawie NOR-145 jego wbudowane
        odczyty głośności i ostrości będą służyć jako sanity-check
        względem naszego potoku \texttt{mosqito}.
\end{itemize}

\end{document}
